\section{ECCOR}
\label{sECCOR}

\hypertarget{sECCORhy}{The}
ECCOR\index{ECCOR|textbf} module of NJOY is used to prepare libraries for the
reactor-physics code ECCO\cite{ecco}.\index{ECCO} ECCO is a lattice code
developed at the Commissariat \`a l'\'Energie Atomique (CEA) and is dedicated
to the analysis of fast reactor lattices.

The ECCOR module creates an ECCO format library named {\sl eccolib} from the
GENDF files output by the CALENR module. It is assumed that an input GENDF file
contains all neutron cross sections, scatter matrices and CALENDF probability
tables for a single isotope/element to be placed on an ECCO library. The code
reads data from the GENDF tape, including probability table information, and
restructures it to the ECCO format. It also forms cross-sections which are
executes final consistency checks and corrects some data which are slightly
inconsistent\cite{eccolib}.

\subsection{User Input}
\label{ssECCOR_UserInp}

The following description of the user input is reproduced from
the comment cards at the beginning of the ECCOR module.
\index{ECCOR!ECCOR input}
\index{input!ECCOR}

\small
\begin{ccode}
  !---input specifications (free format)----
  !
  ! card 1 file units
  !   ning      input gendf unit (>0 formatted / <0 unformatted)
  !   ninr      input unit of reference table (=0 for create option)
  !   nine      input unit of eccolib (=0 for create option)
  !   nour      output unit of reference table
  !   nout      output unit of eccolib
  !   iprint    long print option (0/1=minimum/maximum) (default=0)
  !
  ! card 2 eccolib structure
  !   lrecl     record length (bytes) for library (usually 4096)
  !   lwordi    number of bytes in an integer word (usually 4)
  !   lwordc    cumber of characters in a string (usually 16)
  !   lwordr    number of bytes in a real(kr) word (usually 4)
  !
  ! card 3 group block structure (2 cards)
  !   nreci     number of group blocks
  !   ing       array of integer, number of groups stored together in each
  !             group block. The sum of nreci values must equal the total
  !             number of energy groups in the eccolib.
  !
  ! card 4 title
  !   title     char(len=80) description of this library
  !
  ! card 5 data source for element
  !   sysour    char(len=16), data source for this element
  !
  ! card 6 material number
  !   matno     material number to be processed from GENDF tape
  !
  ! card 7 material number
  !   type      type of nuclide (either 'FERTILE', 'FISSILE' or
  !             'NON-FISSILE')
  !
  ! card 8 reaction number of thermal scattering matrices (2 cards)
  !   mtherm    number of MT indices
  !   ntherm    mtherm values of MT incices (221 is the value for free
  !             gas scattering)
  !
  ! card 9 isotope card
  !   iel       identifier number for this isotope
  !   awt       atomic weight
  !   sycor     char(len=16) isotope name
  !   sybel     char(len=16) element name
  !
  ! card 10 isotope properties
  !   ecng      non-gamma capture energy (mev)
  !   ecg       gamma capture energy (mev)
  !   efng      non-gamma fission energy (mev)
  !   efg       gamma fission energy (Mev)
  !   dis       disintegration constant (s-1)
  !
  ! card11 vector cross section names (2 cards)
  !   nsmic2    number of x-section type names. If nsmic2=0, use 'TOTAL',
  !             'TRANSPORT','CAPTURE','ELASTIC','INELASTIC','FISSION',
  !             'NU*FISSION','N,XN','AXIAL_TRANSPORT','RADIAL_TRANSPORT'
  !   listmi    user-defined char(len=16) names (if nsmic2 > 0)
  !
  ! card12 matrix cross section names (2 cards)
  !   nsmtx     number of x-section type names. If nsmtx=0, use 'ELASTIC',
  !             'INELASTIC','N,XN'
  !   listmx    user-defined char(len=16) names (if nsmtx > 0)
  !
  ! card13 response function names (2 cards)
  !   nfr       number of response function names. If nsmtx=0, use 'N,P',
  !             'N,ALPHA','N,2N','N,3N','N,4N','N,GAMMA','N,D','N,T',
  !             'N,HE3','N,2ALPHA','N,3ALPHA','N,2P',N,P+ALPHA',
  !             'N,T+2ALPHA','N,D+2ALPHA','N,F','N,N+F','N,2N+F','N,3N+F',
  !             'N,N+ALPHA','N,N+3ALPHA','N,2N+ALPHA','N,3N+ALPHA',
  !             'N,N+P','N,N+2ALPHA','N,2N+2ALPHA','N,N+D','N,N+T',
  !             'N,N+HE3','N,N+D+2ALPHA','N,N+T+2ALPHA','N,2NP','N,3NP',
  !             'N,PD','N,PT,’N,ANYTHING’
  !   listfr    user-defined char(len=16) names (if nfr > 0)
  !
  ! card14 response function names (2 cards)
  !   nsgp      number of gamma production x-section names
  !   listgp    user-defined char(len=16) names (if nsgp > 0)
  !
  ! card15 response function names
  !   nmpn      maximum order of pn
  !
  ! card16      gamma energy mesh (2 cards)
  !   nfgg      number of gamma energy groups
  !   eg        energy limits (nfgg+1 values)
  !  ...
  !
  !            cards 9 to 16 are set for the first isotope in the library.
  !            Repeat cards 9 and 10 for all other isotopes present in the
  !            library.
  !
  !-----------------------------------------------------------------

\end{ccode}
\normalsize

The following sample run created an eccolib on tape41 and includes data for
Pu239. The numbers on the left are for reference; they are not part of the
input.

\small
\begin{ccode}
  1.   --
  2.   -- Pu239
  3.   --
  4.   moder
  5.    40 -21
  6.   reconr
  7.    -21 -22
  8.    'pendf tape from /tmp/endfb8r0'/
  9.    9437 1/
 10.    0.001  0.  0.005/
 11.    'Pu239 from /tmp/endfb8r0 at Thu Mar 26 16:01:52 2020' /
 12.    0/
 13.   broadr
 14.    -21 -22 -23
 15.    9437 2/
 16.    0.001/
 17.    2.930000E+02 5.500000E+02/
 18.    0/
 19.   purr
 20.    -21 -23 -24
 21.    9437 2 1 20 32/
 22.    2.930000E+02 5.500000E+02/
 23.    1.000000E+10/
 24.    0/
 25.   thermr
 26.    0 -24 -35
 27.    0 9437 16 2 1 0 0 1 221 0
 28.    2.930000E+02 5.500000E+02/
 29.    0.001 4.0
 30.   groupr
 31.    -21 -35 0 -26 0 0
 32.    9437 35 0 -4 1 2 1 1 0
 33.    'Pu239 from /tmp/endfb8r0 at Thu Mar 26 16:06:27 2020' /
 34.    2.930000E+02 5.500000E+02/
 35.    1.000000E+10/
 36.    677.2873 10.8897 20000 / Homog. Flux Calc.Param
 37.    0.2 0.0253 820.3e3 1.40e6 / iwt=4 parameters
 38.    3/
 39.    3 221 /
 40.    3 452 /
 41.    3 455 /
 42.    5 455 /
 43.    6 /
 44.    6 221 /
 45.    0/
 46.    3/
 47.    3 221 /
 48.    3 452 /
 49.    3 455 /
 50.    5 455 /
 51.    6 /
 52.    6 221 /
 53.    0/
 54.    0/
 55.   calenr
 56.    -21 -35 -26 -27 2 1/
 57.    9437 2 /
 58.    2.930000E+02 5.500000E+02 /
 59.    2.767920E+00 1.227734E+05 /
 60.    0/
 61.   moder
 62.    -27 51
 63.   eccor
 64.    -27 0 0 40 41 1 /
 65.    4096 4 16 4 /
 66.    73 /
 67.      7   7   8   8   8   8   8   9  10   9   9   9
 68.      9   9   9   9  11  16  17  17  18  18  18  18
 69.     19  19  19  19  20  20  20  20  20  21  21  21
 70.     21  22  22  22  23  23  23  24  24  25  25  26
 71.     26  27  28  28  29  30  31  32  34  35  36  38
 72.     40  43  43  48  55  63  74  90 144  69  15  67
 73.     49 /
 74.    'FINE GROUP ENDF-B-VIII.1 LIBRARY' /
 75.    'Pu239 data' /
 76.    9437 /
 77.    'FISSILE' /
 78.    1 /
 79.    221 /
 80.    9437 239.052 'Pu239' 'Pu' /
 81.    0.0 6.5336 189.72 11.910 9.1104E-13 /
 82.    0 /
 83.    0 /
 84.    0 /
 85.    0 /
 86.    1 /
 87.    0 /
 88.   stop
\end{ccode}
\normalsize

Here is the description of the input file necessary to add a nuclide to an
existing eccolib.

\small
\begin{ccode}
  1.   eccor
  2.    -27 38 39 40 41 1 /
  3.    4096 4 16 4 /
  4.    73 /
  5.      7   7   8   8   8   8   8   9  10   9   9   9
  6.      9   9   9   9  11  16  17  17  18  18  18  18
  7.     19  19  19  19  20  20  20  20  20  21  21  21
  8.     21  22  22  22  23  23  23  24  24  25  25  26
  9.     26  27  28  28  29  30  31  32  34  35  36  38
 10.     40  43  43  48  55  63  74  90 144  69  15  67
 11.     49 /
 12.    'Pu240 data' /
 13.    9440 /
 14.    'FISSILE' /
 15.    1 /
 16.    221 /
 17.    9440 240.05381 Pu240 Pu /
 18.    0.0 5.24152 186.80 12.67 3.3477446E-12 /
 19.   stop
\end{ccode}
\normalsize

\subsection{Coding Details}
\label{ssECCOR_details}

The main entry point is subroutine \cword{eccor} exported by
module \cword{eccom}\index{modules!eccom@{\ty eccom}}.
The code begins by reading the user's input. It then locates the
position for the new material on the GENDF\index{GENDF}
tape (if any) and copies them to the eccolib, defined as a binary
direct access file.

To create a new eccolib, ECCOR should first be run in creation
mode.
This option reads
\begin{itemize}
\item the information required for the reference table package:
\begin{itemize}
  \item primary cross section types
  \item response function types
  \item matrix types
  \item sub-group types
  \item fine group energy boundaries
\end{itemize}
\item data describing the first nuclide to be included in the library.
\item a GENDF tape containing the nuclide to be included in the library,
including probability table information produced by module CALENR.
\end{itemize}

Once an eccolib has been produced for one nuclide ECCOR should then be
run in addition mode. This option reads
\begin{itemize}
  \item data describing the nuclide to be added
  \item a GENDF tape containing the nuclide to be added
  \item an existing eccolib
  \item reference table package describing the input eccolib and outputs
\end{itemize}

At completion, ECCOR outputs
\begin{itemize}
  \item a new eccolib
  \item a reference table package describing the eccolib.
\end{itemize}

\cleardoublepage
